% Generated draft: requires mathematical and source review.
\documentclass[11pt]{article}
\usepackage[margin=25mm]{geometry}
\usepackage{fontspec}
\setmainfont{Latin Modern Roman}
\usepackage{amsmath,amssymb,mathtools,unicode-math}
\setmathfont{Latin Modern Math}
\usepackage{enumitem,needspace,xurl,hyperref}
\setlength{\parindent}{0pt}
\setlength{\parskip}{5pt}
\setlength{\emergencystretch}{3em}
\newcommand{\N}{\mathbb{N}}
\newcommand{\Z}{\mathbb{Z}}
\newcommand{\Q}{\mathbb{Q}}
\newcommand{\R}{\mathbb{R}}
\newcommand{\C}{\mathbb{C}}
\newcommand{\F}{\mathbb{F}}
\newcommand{\A}{\mathbb{A}}
\newcommand{\PP}{\mathbb{P}}
\newcommand{\E}{\mathbb{E}}
\newcommand{\eps}{\varepsilon}
\newcommand{\dd}{\,\mathrm{d}}
\begin{document}
\section*{UnsolvedMath ID: 30004372}
\textbf{Source problem code:} catalogue entry
\subsection{Definitions and mathematical statement}
An algebra is hereditary if, in the chosen convention for its modules, every submodule of a projective module is projective. The source does not specify left versus right modules, the base field or division algebra, or finiteness assumptions on the algebra and its modules.

For an algebra \(A\), an \(A\)-module \(M\), and a suitable dimension datum \(\mathbf e\), the submodule quiver Grassmannian \(\operatorname{Gr}^{A}_{\mathbf e}(M)\) parametrizes submodules \(U\subseteq M\) of dimension datum \(\mathbf e\). This notation is conditional on specifying the base, the dimension convention, and the geometric meaning of the parameter space; the supplied record does not provide those choices for general hereditary algebras.

\textbf{Problem (source formulation).} For a general hereditary algebra \(A\), determine an appropriate analogue of the Euler characteristic of \(\operatorname{Gr}^{A}_{\mathbf e}(M)\) that can be used to define the Caldero--Chapoton map in all relevant cases. The supplied assessment explicitly includes real and quaternionic cases among the unresolved extensions and reports that a polynomial Euler characteristic analogue is available over finite fields.

\textbf{Unresolved scope.} The phrase ``all relevant cases'' is not defined. In particular, the source supplies neither a precise class of algebras and modules nor a target ring, normalization, compatibility axioms, or formula for the desired Caldero--Chapoton map. Consequently, this is a request to construct a suitable invariant and map, rather than a fully specified existence conjecture with fixed hypotheses. The reported finite-field result and unresolved extensions are supplied literature qualifications, not claims independently verified in this draft.
\subsection{Short English statement}
Find a replacement for the Euler characteristic of spaces of submodules that allows the Caldero--Chapoton map to be defined for general hereditary algebras. The supplied record reports a finite-field polynomial analogue and leaves extensions, including real and quaternionic cases, unresolved. The intended class of algebras, geometric conventions, and required properties of the replacement remain unspecified.
\subsection{Sources}
\begin{itemize}
\item Local UnsolvedMath registry, permanent ID \(30004372\), title ``Caldero--Chapoton Maps for General Hereditary Algebras,'' source problem code \texttt{catalogue entry}, in the collection ``Oberwolfach Reports Open Problems.'' No individual proposer is identified; the supplied proposed year is 2020.
\item Supplied report citation: \emph{Representation Theory of Quivers and Finite Dimensional Algebras} (2021), \emph{Oberwolfach Reports}, pp.~143--230, DOI \texttt{10.4171/owr/2020/3}. Supplied source URL and literature-source link, labelled ``Oberwolfach report'': \url{https://doi.org/10.4171/owr/2020/3}. The supplied background describes Oberwolfach Reports as documenting mathematical workshops and their open-problem discussions.
\item Literature assessment checked at \texttt{2026-08-22}: the supplied curation labels the problem open, reports a polynomial Euler characteristic analogue over finite fields, and leaves extension to other fields, including real and quaternionic cases, unresolved. This is supplied dated curation, not an independently verified assessment. The background explicitly requests independent verification of time-sensitive status and cited claims.
\item The registry labels the statement \texttt{corrected\_verified} and reports manual semantic, typographic, and LaTeX review. These are supplied editorial qualifications; this draft does not independently verify the report, its mathematical claims, or the current problem status.
\end{itemize}
\end{document}
